\documentclass[11pt,a4paper]{article}
\newcommand{\fichetype}{Corrigé détaillé}
% Thème A4 commun, librement modifiable.
\usepackage[T1]{fontenc}
\usepackage[utf8]{inputenc}
\usepackage[french]{babel}
\usepackage{lmodern,amsmath,amssymb,mathtools,tikz,pgfplots}
\usepackage[margin=2.2cm,headheight=15pt]{geometry}
\usepackage{xcolor,fancyhdr,enumitem,hyperref,needspace}
\usetikzlibrary{arrows.meta,positioning,calc}
\pgfplotsset{compat=1.18,width=11cm,height=6cm}
\definecolor{Ink}{HTML}{202124}
\definecolor{Blue}{HTML}{356B9E}
\definecolor{Muted}{HTML}{737579}
\color{Ink}
\renewcommand{\familydefault}{\sfdefault}
\setlength{\parindent}{0pt}
\setlength{\parskip}{6pt}
\setlist{itemsep=3pt,topsep=4pt}
\pagestyle{fancy}\fancyhf{}
\fancyhead[L]{\small\color{Muted}MATHS LYCÉE}
\fancyhead[R]{\small\color{Muted}2026--2027}
\fancyfoot[L]{\small\color{Muted}\fichetype}
\fancyfoot[R]{\small\thepage}
\renewcommand{\headrulewidth}{0.3pt}
\hypersetup{colorlinks=true,urlcolor=Blue,linkcolor=Blue}
\newcommand{\R}{\mathbb{R}}
\newcommand{\N}{\mathbb{N}}
\newcommand{\E}{\mathbb{E}}
\newcommand{\V}{\mathrm{V}}
\newcommand{\exercice}[3]{\Needspace{6\baselineskip}\vspace{0.5cm}{\color{Blue}\small\textbf{EXERCICE #1}\hfill #2}\par{\large\bfseries #3}\par}

\begin{document}
{\small\color{Blue}PREMIÈRE SPÉCIALITÉ}\par{\LARGE\bfseries Suites arithmétiques et géométriques}\par{\large Corrigé détaillé}\par\vspace{0.2cm}\hrule\vspace{0.4cm}
\par\noindent\begin{minipage}{\linewidth}\exercice{1}{Application}{Une suite arithmétique}
{\small\color{Muted}\textbf{Énoncé.} $u_0=12$ et $r=-3$. Calculer $u_7$, écrire $u_n$ et déterminer le premier rang tel que $u_n<0$.}
\par\medskip
\textbf{Résolution.}\par
La formule est $u_n=12-3n$, donc $u_7=-9$. L’inégalité $12-3n<0$ donne $n>4$. Le premier rang entier est $5$ : $u_4=0$ ne satisfait pas l’inégalité stricte et $u_5=-3<0$.
\end{minipage}\par
\par\noindent\begin{minipage}{\linewidth}\exercice{2}{Entraînement}{Retrouver la raison}
{\small\color{Muted}\textbf{Énoncé.} Une suite arithmétique vérifie $u_3=11$ et $u_8=31$. Trouver $r$, $u_0$ et $u_{20}$.}
\par\medskip
\textbf{Résolution.}\par
Entre les rangs $3$ et $8$, il y a cinq accroissements : $31-11=5r$, donc $r=4$. Puis $u_0=u_3-3r=11-12=-1$. Enfin $u_{20}=-1+20\times4=79$.
\end{minipage}\par
\par\noindent\begin{minipage}{\linewidth}\exercice{3}{Application}{Sommes arithmétiques}
{\small\color{Muted}\textbf{Énoncé.} Calculer $1+\cdots+80$, puis $7+11+15+\cdots+63$.}
\par\medskip
\textbf{Résolution.}\par
La première somme vaut $80\times81/2=3240$. Pour la seconde, le terme de rang $k$ à partir de $0$ est $7+4k$. Résoudre $7+4k=63$ donne $k=14$, soit quinze termes. La somme vaut $15(7+63)/2=525$.
\end{minipage}\par
\par\noindent\begin{minipage}{\linewidth}\exercice{4}{Application}{Une suite géométrique}
{\small\color{Muted}\textbf{Énoncé.} $v_1=6$, $q=1/2$. Calculer $v_5$, puis $v_1+\cdots+v_5$.}
\par\medskip
\textbf{Résolution.}\par
Il faut quatre multiplications pour passer de $1$ à $5$ : $v_5=6(1/2)^4=3/8$. La somme comporte cinq termes : $6[1-(1/2)^5]/(1-1/2)=12\times31/32=93/8=11,625$.
\end{minipage}\par
\par\noindent\begin{minipage}{\linewidth}\exercice{5}{Entraînement}{Variations comparées}
{\small\color{Muted}\textbf{Énoncé.} Étudier les variations de $a_n=4(0,7)^n$, $b_n=-3\times2^n$ et $c_n=5(-1/2)^n$.}
\par\medskip
\textbf{Résolution.}\par
$a_n>0$ et sa raison est comprise entre $0$ et $1$ : elle décroît strictement. Pour $b$, la différence vaut $-3\times2^n<0$, donc elle décroît aussi, malgré une raison supérieure à $1$. Pour $c$, les premiers termes sont $5,-2,5,1,25$ : ils descendent puis montent ; l’alternance se poursuit, donc la suite n’est pas monotone.
\end{minipage}\par
\par\noindent\begin{minipage}{\linewidth}\exercice{6}{Entraînement}{Placement et versement}
{\small\color{Muted}\textbf{Énoncé.} Un capital fictif de $1500$ euros augmente de $4\,\%$ par an sans versement. Donner sa valeur après trois ans et le taux total d’évolution. Comparer à trois hausses additionnées.}
\par\medskip
\textbf{Résolution.}\par
La suite est géométrique de raison $1,04$. Après trois ans : $C_3=1500(1,04)^3=1687,296$, soit $1687,30$ euros au centime. Le coefficient total est $1,124864$, donc le taux vaut $12,4864\,\%$. Additionner les taux donnerait $12\,\%$ : cela oublie que les hausses portent sur des montants successivement plus grands.
\end{minipage}\par
\par\noindent\begin{minipage}{\linewidth}\exercice{7}{Entraînement}{Algorithme de cumul}
{\small\color{Muted}\textbf{Énoncé.} On initialise $v=2$, $S=0$. On répète six fois : $S\leftarrow S+v$, puis $v\leftarrow3v$. Quelle somme calcule-t-on ? La calculer sans boucle.}
\par\medskip
\textbf{Résolution.}\par
La première addition porte sur $v_0=2$. Les six additions portent donc sur $v_0$ à $v_5$, de raison $3$. La somme est $2(1-3^6)/(1-3)=3^6-1=728$. À la fin, la variable $v$ contient déjà $v_6$, mais ce dernier terme n’a pas été ajouté.
\end{minipage}\par
\par\noindent\begin{minipage}{\linewidth}\exercice{8}{Approfondissement}{Remboursement : changement de suite}
{\small\color{Muted}\textbf{Énoncé.} Une dette fictive suit $D_{n+1}=1,1D_n-100$, avec $D_0=500$. Poser $v_n=D_n-1000$. Montrer que $v$ est géométrique et exprimer $D_n$, tant que le modèle décrit une dette.}
\par\medskip
\textbf{Résolution.}\par
$v_{n+1}=D_{n+1}-1000=1,1D_n-1100=1,1(D_n-1000)=1,1v_n$. La raison est $1,1$ et $v_0=-500$. Ainsi $v_n=-500(1,1)^n$ et $D_n=1000-500(1,1)^n$. La dette décroît car $D_{n+1}-D_n=-50(1,1)^n<0$. Quand la formule devient négative, il faut adapter le dernier remboursement : une dette remboursée ne se poursuit pas en montant négatif dans cette situation.
\end{minipage}\par
\end{document}
